\documentclass[11pt]{article}

% ============ FONTS ============
\usepackage{fontspec}
\setmainfont{SpaceGrotesk-Regular.otf}[
  Path=./fonts/,
  BoldFont=SpaceGrotesk-Bold.otf,
  FontFace={lt}{n}{SpaceGrotesk-Light.otf},
  FontFace={mb}{n}{SpaceGrotesk-Medium.otf}]
\newfontfamily\grotesklight{SpaceGrotesk-Light.otf}[Path=./fonts/]
\newfontfamily\groteskmed{SpaceGrotesk-Medium.otf}[Path=./fonts/]
\setmonofont{IBMPlexMono-Regular.otf}[
  Path=./fonts/,
  BoldFont=IBMPlexMono-Bold.otf,
  FontFace={md}{n}{IBMPlexMono-Medium.otf}]

% ============ GEOMETRY & COLOR ============
\usepackage[a4paper,margin=2.3cm,top=2.5cm,bottom=2.5cm]{geometry}
\usepackage[dvipsnames,table]{xcolor}

% Observatory palette (dark slate)
\definecolor{obsbg}{HTML}{0F172A}       % deep slate (title page)
\definecolor{obscard}{HTML}{1E293B}     % card slate
\definecolor{obsink}{HTML}{1A2333}      % near-black ink for body text
\definecolor{obsmuted}{HTML}{64748B}    % muted slate
\definecolor{obsline}{HTML}{CBD5E1}     % light rule
\definecolor{obsaccent}{HTML}{38BDF8}   % sky accent
\definecolor{obsviolet}{HTML}{818CF8}   % indigo accent
\definecolor{obspaper}{HTML}{F8FAFC}    % paper

% Confidence chip colors
\definecolor{cbelegt}{HTML}{059669}     % emerald
\definecolor{cplausibel}{HTML}{0284C7}  % blue
\definecolor{cumstritten}{HTML}{D97706} % amber
\definecolor{cspekulativ}{HTML}{9333EA} % purple

% ============ PACKAGES ============
\usepackage{fancyhdr}
\usepackage{titlesec}
\usepackage{enumitem}
\usepackage{tcolorbox}
\tcbuselibrary{skins,breakable}
\usepackage{amssymb}
\usepackage{booktabs}
\usepackage{array}
\usepackage{longtable}
\usepackage{ragged2e}
\usepackage{microtype}
\usepackage[autostyle=false,german=quotes]{csquotes}
\usepackage{hyperref}
\hypersetup{colorlinks=true,linkcolor=obsaccent,urlcolor=obsviolet,citecolor=obsaccent}


\color{obsink}

% ============ HEADERS ============
\pagestyle{fancy}
\fancyhf{}
\renewcommand{\headrulewidth}{0.4pt}
\renewcommand{\headrule}{\hbox to\headwidth{\color{obsline}\leaders\hrule height \headrulewidth\hfill}}
\fancyhead[L]{\footnotesize\color{obsmuted}\scshape Die Frau vor dem Recht}
\fancyhead[R]{\footnotesize\color{obsmuted}\scshape Zusatzstudie · Die Funktionskritik}
\fancyfoot[C]{\footnotesize\color{obsmuted}\thepage}

% ============ SECTION STYLES ============
\titleformat{\section}
  {\groteskmed\Large\color{obsink}}
  {\color{obsaccent}\S\,\thesection}{0.7em}{}
  [\vspace{-0.6em}{\color{obsline}\rule{\textwidth}{0.4pt}}]
\titleformat{\subsection}
  {\groteskmed\large\color{obsink}}
  {\color{obsviolet}\thesubsection}{0.6em}{}
\titlespacing{\section}{0pt}{1.4em}{0.8em}

% ============ CONFIDENCE CHIPS ============
\newtcbox{\chip}[1][obsmuted]{on line,boxsep=0pt,left=5pt,right=5pt,top=1.5pt,bottom=1.5pt,
  colframe=#1,colback=#1,coltext=white,boxrule=0pt,arc=2pt,fontupper=\ttfamily\scriptsize\bfseries}
\newcommand{\belegt}{\chip[cbelegt]{belegt}}
\newcommand{\plausibel}{\chip[cplausibel]{plausibel}}
\newcommand{\umstritten}{\chip[cumstritten]{umstritten}}
\newcommand{\spekulativ}{\chip[cspekulativ]{spekulativ}}

% ============ CALLOUT BOXES ============
\newtcolorbox{gegen}{
  enhanced,breakable,
  colback=obspaper,colframe=cumstritten,
  boxrule=0pt,leftrule=3pt,arc=1pt,
  left=12pt,right=12pt,top=8pt,bottom=8pt,
  fontupper=\normalsize,
  overlay unbroken and first={
    \node[anchor=north west,font=\ttfamily\scriptsize\bfseries\color{cumstritten}]
    at ([xshift=12pt,yshift=-4pt]frame.north west){GEGENDARSTELLUNG};},
  before upper={\vspace{1.1em}}}

\newtcolorbox{merksatz}{
  enhanced,breakable,
  colback=obscard,colframe=obscard,
  coltext=obspaper,
  boxrule=0pt,arc=3pt,
  left=16pt,right=16pt,top=12pt,bottom=12pt,
  fontupper=\groteskmed\large}

\newtcolorbox{datenbox}[1][]{
  enhanced,breakable,
  colback=obspaper,colframe=obsaccent,
  boxrule=0pt,leftrule=3pt,arc=1pt,
  left=12pt,right=12pt,top=8pt,bottom=8pt,
  coltitle=obsaccent,fonttitle=\ttfamily\scriptsize\bfseries,
  title={#1},attach title to upper=\\}

% ============ TABLE HELPERS ============
\newcolumntype{L}[1]{>{\RaggedRight\arraybackslash}p{#1}}
\renewcommand{\arraystretch}{1.35}

\setlist[itemize]{leftmargin=1.3em,itemsep=2pt,topsep=3pt}
\setlength{\parindent}{0pt}
\setlength{\parskip}{0.55em}

\begin{document}

% ================= TITLE PAGE =================
\begin{titlepage}
\thispagestyle{empty}
\newgeometry{margin=0pt}
\colorbox{obsbg}{%
\begin{minipage}[t][\paperheight][t]{\paperwidth}
\color{obspaper}
\vspace{2.6cm}
\hspace{2.3cm}\begin{minipage}{0.82\textwidth}
{\ttfamily\footnotesize\color{obsaccent}OBSERVATORY \quad·\quad ZUSATZSTUDIE ZU BAND VII \quad·\quad DIAGNOSE UND PROGRAMM}\\[0.3em]
{\color{obsmuted}\rule{\linewidth}{0.4pt}}

\vspace{2.4cm}
{\grotesklight\fontsize{15}{18}\selectfont\color{obsline} Die Frau vor dem Recht --- Werkstatt}\\[0.7em]
{\groteskmed\fontsize{40}{46}\selectfont Die\\[0.15em] Funktionskritik}\\[1.1em]
{\grotesklight\fontsize{16}{22}\selectfont\color{obsline}\itshape
Entwurf einer Disziplin. Korrektur, Erweiterung\\ oder Systemkritik des Feminismus?}

\vspace{2.2cm}
{\color{obsmuted}\rule{\linewidth}{0.4pt}}\\[0.5em]
{\grotesklight\normalsize\color{obsline}
Zwei Zusatzdokumente haben das Verhältnis der Funktionslogik-These zum
intersektionalen Feminismus vermessen: echte Zieldeckungen, echte Kollisionen,
und eine Leerstelle, die sich als \textit{Feature mit Bug-Anteilen} erwies.
Diese Programmschrift zieht die Konsequenz und entwirft die fehlende
Disziplin --- ihren Gegenstand, ihre Axiome, ihre Methode, ihre Grenzen und
ihre stärksten Einwände gegen sich selbst.}

\vspace{1.6cm}
{\ttfamily\scriptsize\color{obsline}
KONFIDENZMARKER\\[0.5em]
\textcolor{cbelegt}{$\blacksquare$}~belegt · quellennah/empirisch gesichert \qquad
\textcolor{cplausibel}{$\blacksquare$}~plausibel · begr\"undete Deutung\\[0.3em]
\textcolor{cumstritten}{$\blacksquare$}~umstritten · aktive Forschungskontroverse \qquad
\textcolor{cspekulativ}{$\blacksquare$}~spekulativ · Hypothese jenseits der Evidenz}

\vfill
{\ttfamily\footnotesize\color{obsmuted}6.~Juli~2026 \quad·\quad Leitformel: Rechtsstellung folgt Funktionsstellung}
\vspace{2cm}
\end{minipage}
\end{minipage}}
\restoregeometry
\end{titlepage}

\tableofcontents
\thispagestyle{empty}
\clearpage

\input{fk_body.tex}

\end{document}
